%% Appendix H: the signal task redrawn to show how reasoning is scored (2026-10-26).
%% Top half as in signal_game.tex. Bottom: the solver answers ONE action; the rules consistent with the twenty clues
%% give the share s of each action; the reasoning score is 1 when the answer is the top-share action.
%% The scored truth is one consistent rule drawn at random, so each action is the truth with probability s.
%% Same example as signal_game.tex: s = (0, .70, 0, .30), the solver answered go_right (the top action).
\begingroup
\definecolor{sgyel}{HTML}{D9A400}
\newcommand{\sgsig}[4]{%
  \begin{scope}[shift={#1}]
    \foreach \j in {1,...,#4}{%
    \begin{scope}[xshift={.3*(\j-(#4+1)/2)*1cm}]
    \ifthenelse{\equal{#3}{circle}}{\fill[#2] (0,0) circle (.115);}{}%
    \ifthenelse{\equal{#3}{square}}{\fill[#2] (-.105,-.105) rectangle (.105,.105);}{}%
    \ifthenelse{\equal{#3}{triangle}}{\fill[#2] (90:.145)--(210:.145)--(330:.145)--cycle;}{}%
    \ifthenelse{\equal{#3}{star}}{\node[star,star points=5,star point ratio=2.3,fill=#2,inner sep=0pt,minimum size=.28cm] at (0,0) {};}{}%
    \end{scope}}
  \end{scope}}
\newcommand{\sgact}[2]{%
  \begin{scope}[shift={#1}]
    \ifthenelse{\equal{#2}{L}}{\draw[black!70,line width=.9pt,-{Stealth[length=3pt]}] (.14,0)--(-.16,0);}{}%
    \ifthenelse{\equal{#2}{R}}{\draw[black!70,line width=.9pt,-{Stealth[length=3pt]}] (-.14,0)--(.16,0);}{}%
    \ifthenelse{\equal{#2}{S}}{\fill[black!70] (-.07,-.07) rectangle (.07,.07);}{}%
    \ifthenelse{\equal{#2}{J}}{\draw[black!70,line width=.9pt,-{Stealth[length=3pt]}] (0,-.13)--(0,.14);}{}%
  \end{scope}}
\newcommand{\sgex}[5]{%
  \begin{scope}[shift={#1}]
    \draw[black!25,fill=white,rounded corners=1.5pt] (.1,-.18) rectangle (2.55,.18);
    \draw[sgorange!40,fill=sgorange!12,rounded corners=1.5pt] (2.72,-.18) rectangle (3.55,.18);
    \sgsig{(1.325,0)}{#2}{#3}{#4}
    \sgact{(3.135,0)}{#5}
  \end{scope}}
\newcommand{\sgrule}[2]{%
  \begin{scope}[shift={#1},rotate=#2]
    \draw[black!35,fill=white,rounded corners=1pt,line width=.4pt] (-.2,-.17) rectangle (.2,.17);
    \foreach \y in {.093,.031,-.031}{\fill[black!25] (-.15,\y-.018) rectangle (.02,\y+.018);
      \fill[sgorange!45] (.06,\y-.018) rectangle (.15,\y+.018);}
    \fill[sgorange!45] (.06,-.111) rectangle (.15,-.075);
  \end{scope}}
\begin{tikzpicture}[x=1cm,y=1cm]
%% ---- top left: example cards
\node[sgfont=7,text=black!70] at (1.83,4.25) {단서 (20개 중 4개)};
\node[sgfont=6.3,text=black!70] at (1.325,3.93) {시그널};
\node[sgfont=6.3,text=sgorange!85!black] at (3.135,3.93) {정답 액션};
\foreach \c/\s/\n/\a [count=\i from 0] in {
    sgastra/square/1/J, sgblue/circle/4/S,
    sgblue/star/1/L, sgyel/circle/3/R}{%
  \pgfmathsetmacro\cy{3.57-\i*0.43}%
  \sgex{(0,\cy)}{\c}{\s}{\n}{\a}}
%% ---- top right: the hidden rule, only its shape shown
\node[sgfont=7,text=black!70] at (6.23,4.25) {숨은 규칙 (모양만 공개)};
\begin{scope}[shift={(4.25,2.06)}]
  \draw[black!30,fill=black!3,rounded corners=3pt] (0,0) rectangle (3.95,1.98);
  \node[sgfont=6.2,text=black!55] at (1.7,1.72) {조건 = 색 · 모양 · 수};
  \node[sgfont=6.2,text=sgorange!85!black] at (3.16,1.72) {액션};
  \foreach \k/\y in {if/1.36,elif/1.0,elif/.64,else/.28}{%
    \node[sgfont=7,anchor=west,font=\ttfamily\fontsize{7}{8}\selectfont,text=sgblue!80!black] at (.18,\y) {\k};
    \ifthenelse{\equal{\k}{else}}{}{\fill[black!22,rounded corners=1pt] (1.05,\y-.11) rectangle (2.35,\y+.11);}
    \fill[sgorange!35,rounded corners=1pt] (2.7,\y-.11) rectangle (3.62,\y+.11);}
\end{scope}
\draw[black!35,line width=.6pt,-{Stealth[length=3.5pt]}] (3.7,3.05)--(4.18,3.05);
%% ---- bottom: the new signal, top centre
\draw[black!30,dashed,line width=.4pt] (0,1.86)--(8.2,1.86);
\node[sgfont=7,text=black!70] at (4.1,1.64) {새 시그널};
\draw[black!45,fill=white,rounded corners=2pt,line width=.8pt] (3.4,.98) rectangle (4.8,1.42);
\sgsig{(4.1,1.2)}{sgblue}{square}{3}
%% ---- bottom left: the solver answers one action
\sgrobot[.62]{(1.57,1.16)}{sgorange}{solver}
\draw[black!40,line width=.6pt,-{Stealth[length=3.5pt]}] (3.32,1.2)--(2.02,1.2);
\node[sgfont=6.3,text=black!60,anchor=north] at (2.67,1.17) {추론};
\draw[sgorange!80,line width=.6pt,-{Stealth[length=3.5pt]}] (1.57,.86)--(1.57,.5);
\draw[sgorange!60,fill=sgorange!12,rounded corners=2pt,line width=.7pt] (.95,-.05) rectangle (2.2,.42);
\node[sgfont=6.6,anchor=west] at (1.02,.19) {답:};
\sgact{(1.75,.19)}{R}
\node[sgfont=6.2,text=black!65] at (1.57,-.28) {행동 하나만 낸다};
%% ---- bottom right: rule candidates, filtered by the twenty clues, give the share s
\node[sgfont=7,text=black!70,anchor=east] at (8.15,1.64) {단서에 맞는 규칙 후보};
\fill[black!4] (6.8,1.12) ellipse (1.3 and .36);
\draw[black!22,line width=.4pt,dash pattern=on 1.5pt off 1pt] (6.8,1.12) ellipse (1.3 and .36);
\foreach \x/\y/\r/\ok in {5.88/1.14/-8/1, 6.34/.98/6/0, 6.8/1.2/-4/1, 7.26/.99/8/1, 7.72/1.15/-6/0}{%
  \sgrule{(\x,\y)}{\r}
  \ifthenelse{\equal{\ok}{1}}{\sgcheck[.55]{(\x+.2,\y+.17)}}{\sgcross[.55]{(\x+.2,\y+.17)}}}
\draw[black!40,line width=.6pt,-{Stealth[length=3.5pt]}] (4.88,1.2)--(5.47,1.1);
\draw[black!60,line width=.6pt,-{Stealth[length=3.5pt]}] (6.62,.75)--(6.62,.52);
\foreach \v/\act/\top [count=\i from 0] in {0/L/0, .70/R/1, 0/S/0, .30/J/0}{%
  \pgfmathsetmacro\bx{5.5+\i*.62}%
  \ifthenelse{\equal{\top}{1}}{\draw[sgblue!80!black,fill=sgblue!25,line width=.7pt] (\bx,-.42) rectangle ++(.38,\v);}%
    {\draw[black!75,line width=.7pt] (\bx,-.42) rectangle ++(.38,\v);}
  \sgact{(\bx+.19,-.56)}{\act}}
\draw[black!55,line width=.5pt] (5.4,-.42)--(7.84,-.42);
\node[sgfont=5.8,text=sgblue!80!black] at (6.31,.4) {최다 행동 (.70)};
\node[sgfont=5.8,text=black!60] at (7.55,0) {.30};
\node[sgfont=6.6] at (6.7,-.84) {규칙들이 고른 비율 $s$};
%% ---- the score in the middle
\draw[black!80,line width=.8pt,-{Stealth[length=4pt]}] (2.3,.19)--(5.3,.19);
\node[sgfont=6.4,text=black!75] at (3.8,.42) {답이 최다 행동인가?};
\node[font=\sffamily\bfseries\fontsize{7}{8}\selectfont,text=black!90] at (3.8,-.06) {예 = 1, 아니오 = 0 (추론 점수)};
\draw[black!20,line width=.4pt] (0,-1.02)--(8.2,-1.02);
\node[sgfont=6.2,text=black!60,anchor=north,align=center] at (4.1,-1.08) {채점용 정답은 단서에 맞는 후보 규칙 하나를 뽑아 정한다.\\ 그래서 각 행동이 정답일 확률 $= s$이고, 최다 행동을 골라도 .70으로만 맞는다.};
\end{tikzpicture}
\endgroup
